%==============================================================================
\section{Getting Datasets --- Lấy bộ dữ liệu cho học máy}
%==============================================================================

\begin{dongco}[title={Mô hình chỉ tốt bằng dữ liệu huấn luyện}]
Mô hình học máy chỉ tốt bằng dữ liệu dùng để huấn luyện.
Vì vậy, có được bộ dữ liệu chất lượng tốt và phù hợp là bước then chốt trong quy trình ML.
Có nhiều kho mã nguồn mở (ví dụ Kaggle) để tải dataset; bạn cũng có thể mua dữ liệu, scrape website hoặc tự thu thập.
Bài này giới thiệu các nguồn dataset và cách lấy chúng.
\end{dongco}

\begin{vidu}[title={Hình minh hoạ dữ liệu cho ML}]
\tsimg{10_du_lieu_trong_ml/img01.jpg}
\end{vidu}

\subsection{Dataset là gì?}

\begin{dinhnghia}[title={Dataset}]
Dataset là tập dữ liệu được tổ chức có cấu trúc, thường dùng để đơn giản hoá phân tích, lưu trữ, xử lý hoặc huấn luyện mô hình ML.
Dataset có thể lưu ở nhiều định dạng: CSV, JSON, file nén, Excel, \ldots
\end{dinhnghia}

\subsection{Các loại dataset}

\begin{phuongphap}[title={Phân loại theo kiểu thông tin}]
\begin{itemize}
  \item \textbf{Tabular}: dữ liệu dạng bảng (hàng/cột).
  \item \textbf{Time series}: dữ liệu theo thời gian (giá cổ phiếu, khí hậu, \ldots).
  \item \textbf{Image}: ảnh cho computer vision (phân loại ảnh, phát hiện đối tượng, segmentation).
  \item \textbf{Text}: văn bản cho NLP (sentiment analysis, text classification, \ldots).
\end{itemize}
\end{phuongphap}

\subsection{Lấy dataset cho Machine Learning}

\begin{ynghia}[title={Tầm quan trọng của bước lấy dữ liệu}]
Lấy dataset là bước rất quan trọng khi xây giải pháp ML.
Dữ liệu là yếu tố bắt buộc để huấn luyện mô hình; chất lượng, số lượng và độ đa dạng của dữ liệu ảnh hưởng mạnh tới hiệu năng mô hình.
\end{ynghia}

\begin{phuongphap}[title={Bốn nguồn chính}]
\begin{itemize}
  \item Open Source Datasets (dataset công khai)
  \item Data Scraping (thu thập tự động từ web/nguồn khác)
  \item Data Purchase (mua dữ liệu)
  \item Data Collection (thu thập thủ công)
\end{itemize}
\end{phuongphap}

\subsection{Dataset công khai / mã nguồn mở phổ biến}

\begin{ynghia}[title={Kho dữ liệu công khai}]
Có nhiều dataset mã nguồn mở dùng cho ML, ví dụ: Kaggle, UCI Machine Learning Repository, Google Dataset Search, AWS Public Datasets.
Các dataset này thường dùng cho nghiên cứu và mở cho cộng đồng.
\end{ynghia}

\begin{phuongphap}[title={Một số nguồn có dữ liệu có cấu trúc, giá trị cao}]
Kaggle Datasets, AWS Datasets, Google Dataset Search, UCI Machine Learning Repository, Microsoft Datasets, Scikit-learn Dataset, HuggingFace Datasets Hub, Government Datasets.
\end{phuongphap}

\subsubsection{Kaggle Datasets}

\begin{ynghia}[title={Kaggle Datasets}]
Kaggle là cộng đồng trực tuyến phổ biến về data science và ML, lưu trữ hơn 23.000 dataset công khai.
Đây là nền tảng được chọn nhiều nhất để lấy dataset vì dễ tìm, tải và đăng dữ liệu.
Kaggle cung cấp dataset tiền xử lý chất lượng cao, phù hợp hầu hết mô hình ML theo nhu cầu người dùng.
Kaggle còn có notebook kèm thuật toán và các loại mô hình pre-trained.
\end{ynghia}

\subsubsection{AWS Datasets}

\begin{ynghia}[title={AWS Datasets}]
Có thể tìm, tải và chia sẻ dataset công khai trong registry open data trên AWS.
Dù truy cập qua AWS, dataset do tổ chức chính phủ, doanh nghiệp và nhà nghiên cứu duy trì, cập nhật.
\end{ynghia}

\subsubsection{Google Dataset Search Engine}

\begin{ynghia}[title={Google Dataset Search}]
Công cụ của Google cho phép tìm dataset từ nhiều nguồn trên web --- công cụ tìm kiếm chuyên cho dataset.
\end{ynghia}

\subsubsection{UCI Machine Learning Repository}

\begin{ynghia}[title={UCI ML Repository}]
Kho dataset của Đại học California, Irvine, dành riêng cho ML.
Hàng trăm dataset từ nhiều miền: time series, classification, regression, recommendation systems, \ldots
\end{ynghia}

\subsubsection{Microsoft Dataset}

\begin{ynghia}[title={Microsoft Dataset}]
Microsoft Research Open Data cung cấp kho dữ liệu trên cloud.
\end{ynghia}

\subsubsection{Scikit-learn Dataset}

\begin{ynghia}[title={Scikit-learn Dataset}]
Scikit-learn cung cấp dataset mẫu (Iris, \ldots) để thử nghiệm.
Dataset mở, dùng học và thử mô hình ML.
\end{ynghia}

\begin{vidu}[title={Nạp dataset Iris từ scikit-learn}]
\begin{lstlisting}
from sklearn.datasets import load_iris
iris = load_iris()
\end{lstlisting}
Đoạn trên nạp dataset Iris vào script Python.
\end{vidu}

\subsubsection{HuggingFace Datasets Hub}

\begin{ynghia}[title={HuggingFace Datasets Hub}]
Hub cung cấp các dataset công khai lớn: ảnh, âm thanh, văn bản, \ldots
\end{ynghia}

\begin{vidu}[title={Cài đặt và nạp dataset với HuggingFace}]
Cài đặt thư viện \texttt{datasets}:
\begin{lstlisting}
pip3 install datasets
\end{lstlisting}
Cú pháp đơn giản:
\begin{lstlisting}
from datasets import load_dataset
ds = load_dataset(dataset_name)
\end{lstlisting}
Ví dụ nạp Iris:
\begin{lstlisting}
from datasets import load_dataset
ds = load_dataset("scikit-learn/iris")
\end{lstlisting}
\end{vidu}

\subsubsection{Government Datasets}

\begin{ynghia}[title={Government Datasets}]
Mỗi quốc gia thường có nguồn dữ liệu chính phủ công khai, thu từ nhiều bộ ngành, nhằm minh bạch và phục vụ nghiên cứu.
\end{ynghia}

\begin{phuongphap}[title={Một số liên kết dataset chính phủ}]
\begin{itemize}
  \item Indian Government Public Datasets
  \item U.S. Government's Open Data
  \item World Bank Data Catalog
  \item European Union Open Data
  \item U.K. Open Data
\end{itemize}
\end{phuongphap}

\subsection{Data Scraping}

\begin{ynghia}[title={Scrape dữ liệu}]
Data scraping là trích xuất tự động dữ liệu từ website hoặc nguồn khác.
Hữu ích khi không có dataset đóng gói sẵn.
Cần đảm bảo scrape \textbf{đạo đức và hợp pháp}, và nguồn đáng tin, chính xác.
\end{ynghia}

\subsection{Data Purchase}

\begin{ynghia}[title={Mua dataset}]
Đôi khi cần mua dataset cho ML.
Nhiều công ty bán dataset theo ngành hoặc use case.
Trước khi mua, cần đánh giá chất lượng và mức độ liên quan tới dự án.
\end{ynghia}

\subsection{Data Collection}

\begin{ynghia}[title={Thu thập thủ công}]
Thu thập thủ công từ nhiều nguồn --- tốn thời gian, cần kế hoạch để dữ liệu chính xác, phù hợp.
Có thể qua khảo sát, phỏng vấn hoặc hình thức thu thập khác.
\end{ynghia}

\subsection{Chiến lược có dataset chất lượng cao}

\begin{phuongphap}[title={Sau khi chọn nguồn, cần đảm bảo chất lượng và liên quan}]
\begin{itemize}
  \item \textbf{Xác định bài toán}: trước khi lấy dữ liệu, làm rõ bài toán ML cần giải --- từ đó biết cần loại dữ liệu gì và lấy ở đâu.
  \item \textbf{Xác định quy mô dataset}: phụ thuộc độ phức tạp bài toán; thường dữ liệu nhiều hơn giúp mô hình tốt hơn, nhưng tránh quá lớn với dữ liệu trùng/không liên quan.
  \item \textbf{Đảm bảo liên quan và chính xác}: dữ liệu từ nguồn tin cậy, đã được kiểm tra.
  \item \textbf{Tiền xử lý}: làm sạch, chuẩn hoá, biến đổi để mô hình dùng hiệu quả.
\end{itemize}
\end{phuongphap}
